%% kapa: exam=1 round=32 year=2021 period=2 form=A subject=appraisal_law subject_ko=감정평가관계법규 seq=36 number=36
%% source: Q-Net 감정평가사 기출문제 (cst00309) artlSeq=5208397 file=extracted/1008/5208397/A형/2021년도 제32회 감정평가사 1차 2교시 A형.pdf
%% conversion: kordoc 4.24.51 -> tools/gen1.py (auto); review=proofread-sonnet; flags=none
\begin{question}{36}
동산·채권 등의 담보에 관한 법령상 담보등기에 관한 설명으로 옳은 것은?
\begin{choices}
  \item 판결에 의한 등기는 등기권리자와 등기의무자가 공동으로 신청하여야 한다.
  \item 등기관이 등기를 마친 경우 그 등기는 등기신청을 접수한 날의 다음 날부터 효력을 발생한다.
  \item 등기관의 결정에 대한 이의신청은 집행정지의 효력이 있다.
  \item 등기관은 자신의 결정 또는 처분에 대한 이의가 이유 없다고 인정하면 3일 이내에 의견서를 붙여 사건을 관할 지방법원에 송부하여야 한다.
  \item 「동산·채권 등의 담보에 관한 법률」에 따른 담보권의 존속기간은 7년을 초과하지 않는 기간으로 이를 갱신할 수 있다.
\end{choices}
\end{question}
\answer{4}
